\section{Cyclotomic elements with compatible divisors and residues}
\label{sec:minus-cyclotomic-elements}

We construct the elements \(X_n\) used in
Section~\ref{sec:minus-derivatives}. Keep its integers \(P,r,b\),
where \(0\le b<P\), and fix a place above \(r\) in
\(\overline{\mathbb Q}\). For a positive integer \(a\) put
\(F_a=\mathbb Q(\mu_a)\), \(G_a=(\mathbb Z/a\mathbb Z)^\times\),
and let \(\sigma_t\) act on roots of unity by \(t\). Write
\([t]_a\) for the least nonnegative residue, and put
\[
   h_a(t)=\left\lfloor b[t]_a/a\right\rfloor-
            \left\lfloor b[-t]_a/a\right\rfloor,\qquad
   \Theta_a=\sum_{t\in G_a} h_a(t)\sigma_t^{-1}.
\]
Use multiplicative group-ring notation \(I^{\Theta_a}\) for ideals
and \(y^{\Theta_a}\) for elements. For a nonzero integral ideal
\(I\) in \(F_a\), seek an element \(X_a(I)\) satisfying
\[
       X_a(I)\overline{X_a(I)}=1,\qquad
       (X_a(I))=I^{\Theta_a}.
\]
Here the bar denotes complex conjugation. Any two choices differ by
a root of unity of order dividing \(2a\). Indeed, their quotient
is an algebraic unit of absolute value one in every embedding,
since complex conjugation commutes with those embeddings.
Kronecker's theorem makes it a root of unity, and the roots of
unity of \(F_a\) have exponent dividing \(2a\).
For \(I=(y)\) we can take \(y^{\Theta_a}\), by
\(h_a(-t)=-h_a(t)\). The two identities are multiplicative in \(I\)
and \(X\), and allow division in the group of fractional ideals.

\subsection{Existence and the Jacobi sum divisor}

Every ideal class of \(F_a\) contains a prime ideal of residue
degree one over \(\mathbb Q\), avoiding any prescribed finite set.
Indeed, the Hilbert class field of \(F_a\) is Galois over
\(\mathbb Q\), since \(F_a/\mathbb Q\) is Galois. Apply Chebotarev
over \(\mathbb Q\) to the Artin element of the class, an element
of the Galois group over \(F_a\). Its restriction to \(F_a\) is
trivial. Choosing the prime above the resulting rational prime to
give that exact Artin element proves the assertion by Hilbert class
field reciprocity. Choose in this manner a degree-one prime
\(\mathfrak p\) in the inverse class of \(I\) over \(p\nmid a\).
If \(I\mathfrak p=(y)\), dividing \(y^{\Theta_a}\) by a choice
at \(\mathfrak p\) gives the desired choice at \(I\).

The reduction map \(\pi:\mathcal O_{F_a}\to\mathbb F_p\) at
\(\mathfrak p\) sends \(\mu_a\) bijectively to the \(a\)-th roots
of unity in \(\mathbb F_p\), so \(a\mid p-1\). Write
\(d=(p-1)/a\), and lift \(x\mapsto x^d\) uniquely to a character
\(\chi:\mathbb F_p^\times\to\mu_a\). Extend every multiplicative
character used in sums by zero at zero. Set
\[
 J_j=\sum_{x\in\mathbb F_p}\chi(x)\chi^j(1-x),\qquad
 J=\prod_{1\le j<b}J_j,\qquad X_a(\mathfrak p)=\overline J/J .
\]
For \(1\le j<b<a\), all of \(\chi,\chi^j,\chi^{j+1}\) are
nontrivial. In any complex embedding the Gauss sums
\(G(\lambda)=\sum_x\lambda(x)e^{2\pi i x/p}\) have
\(\lvert G(\lambda)\rvert^2=p\) for nontrivial \(\lambda\):
put \(x=ty\) in the squared absolute value and sum
\(e^{2\pi i(t-1)y/p}\) over \(y\ne0\), obtaining \(p-1\) at
\(t=1\) and \(-1\) elsewhere. Moreover
\(G(\chi)G(\chi^j)=J_jG(\chi^{j+1})\), by grouping according
to \(x+y\); the terms with \(x+y=0\) sum to zero.
Thus \(J_j\overline J_j=p\). Each \(J_j\) is integral and its
ideal divides the squarefree ideal \((p)\).

For \(t\in G_a\) let \(A=d[t]_a,\ B=d[jt]_a\), both between
\(1\) and \(p-2\). The reduction of \(\sigma_t J_j\) at
\(\mathfrak p\) is \(\sum_x x^A(1-x)^B\).
Expanding the second power and using the power sums of
\(\mathbb F_p^\times\) gives zero if \(A+B<p-1\). Otherwise the
only contributing positive multiple of \(p-1\) is \(p-1\), and
the coefficient is
\(-(-1)^{p-1-A}\binom{B}{p-1-A}\ne0\bmod p\).
The primes \(\sigma_t^{-1}\mathfrak p\) are distinct and exhaust
the primes above \(p\). Hence the exponent of
\(\sigma_t^{-1}\mathfrak p\) in \(J_j\) is
\[
       1-\left\lfloor\big([t]_a+[jt]_a\big)/a\right\rfloor .
\]
The floor is the increment in \(\lfloor j[t]_a/a\rfloor\)
when \(j\) increases to \(j+1\). For \(b\ge1\) telescoping
gives the exponent of \(J\) as \(b-1-\lfloor b[t]_a/a\rfloor\).
Consequently \(\overline J/J\) has the required exponent
\(h_a(t)\) at this prime. At \(b=0\) both the required
exponents and the empty-product exponents are zero. This proves
existence, including the divisor and the archimedean normalization.

\subsection{Norm from a new prime layer}

Fix \(a=Pn\) and an admissible \(q\nmid a\), with \(0\le b<P<q\).
Write \(E=F_a,\ E'=F_{qa},\ u=q\bmod a\). Restriction of the group
ring satisfies
\begin{equation}
       \left.\Theta_{qa}\right|_{G_a}
                  =(1-\sigma_u^{-1})\Theta_a .
 \label{eq:cyclotomic-distribution}
\end{equation}
To see this coefficientwise, the \(q\) representatives lifting
\(t\bmod a\) are \([t]_a+ia\) for \(0\le i<q\). Exactly one,
\(q[u^{-1}t]_a\), is excluded because it is divisible by \(q\).
For \(b>0\), \(ib\bmod q\) permutes the residues modulo \(q\);
summing the floors of
\(b([t]_a+ia)/(qa)\) over all \(i\) therefore gives
\(\lfloor b[t]_a/a\rfloor+\sum_{i<q}\lfloor ib/q\rfloor\).
Indeed the fractional part of \(b[t]_a/a\), divided by \(q\),
does not affect any of these floors, and adding its integral
part permutes the residues in the numerator. Subtract the missing
term \(\lfloor b[u^{-1}t]_a/a\rfloor\), then subtract the same
calculation with \(-t\). The common constant cancels, giving
\eqref{eq:cyclotomic-distribution}. For \(b=0\) it holds identically.

For a nonzero ideal \(I'\) of \(E'\), put \(I=N_{E'/E}I'\).
Norm commutes with conjugation and products of fractional ideals.
Thus the divisor of \(N X_{qa}(I')\) is
\(I^{(1-\sigma_u^{-1})\Theta_a}\). The latter is also the divisor
of \(X_a(I)/\sigma_u^{-1}X_a(I)\). Both elements have product
one with their complex conjugates. Their ratio is therefore an
element of \(\mu_{2a}\), by the uniqueness argument.

\subsection{Simultaneous reduction at every prime over \(q\)}

Suppose \(I'\) is prime to \(q\). We prove, with a single root
\(\eta\in\mu_{2qa}\), the congruence
\[
          X_{qa}(I')\equiv\eta\,\sigma_u^{-1}X_a(I)
                       \pmod{\mathfrak q}
          \qquad(\mathfrak q\mid q\ \text{in }E').
\]
Both sides are units at all such places by their divisors.

First take a degree-one prime \(\mathfrak p'\) above \(p\nmid qa\).
Its relative ideal norm is the contracted prime \(\mathfrak p\).
Use the Jacobi choices already constructed, with characters
\(\chi'\) at level \(qa\) and \(\chi\) at level \(a\), defined by
reduction at \(\mathfrak p'\) and \(\mathfrak p\) respectively.
For \(x\ne0\), \(\chi'(x)^q=\chi(x)\): these are \(a\)-th roots
with equal reductions modulo \(\mathfrak p'\). Also
\((\sigma_u^{-1}\chi(x))^q=\chi(x)\). In the residue field of
any \(\mathfrak q\mid q\), the \(q\)-power map is injective, so
\(\chi'(x)\equiv\sigma_u^{-1}\chi(x)\pmod{\mathfrak q}\).
The same argument applies after complex conjugation. The two pairs
of Jacobi products therefore have equal reductions, term by term.
Each product and its conjugate is a unit at \(q\), since its
factors divide \(p\). Dividing proves the congruence with
\(\eta=1\) for the Jacobi choices. Changing either choice multiplies
the ratio by a single global root of order dividing \(2qa\).

For a principal ideal \((y)\), with \(y\) integral and prime to
\(q\), use the explicit principal generators. Fix
\(\mathfrak q\mid q\) and write \(f(t)\) for the reduction of
\(\sigma_{\tilde t}^{-1}y\) at this place, for any lift
\(\tilde t\in G_{qa}\) of \(t\). This is independent of the lift
because \(\operatorname{Gal}(E'/E)\) fixes residues at \(q\):
\(\mu_q\) reduces to \(1\) and
\(\mathcal O_{E'}=\mathbb Z[\mu_{qa}]\). A lift of \(\sigma_u\)
acts on these residues by \(q\)-power, as this is its action on
the reductions of \(\mu_a\) and \(\mu_q\). Consequently
\(f(ut)^q=f(t)\). All \(f(t)\) are nonzero. Using
\eqref{eq:cyclotomic-distribution}, the reduction of
\(y^{\Theta_{qa}}\) is
\[
 \prod_{t\in G_a} f(t)^{h_a(t)-h_a(u^{-1}t)}
    =\prod_{t\in G_a}\big(f(ut)^{q-1}\big)^{h_a(t)}.
\]
Each of the \(q-1\) conjugates entering \(N_{E'/E}y\) has the
same reduction after applying \(\sigma_{ut}^{-1}\). Hence the
last product is exactly the reduction of
\(\sigma_u^{-1}((Ny)^{\Theta_a})\).
The equality holds at every place above \(q\).

For general \(I'\) prime to \(q\), choose a degree-one
\(\mathfrak p'\) in its inverse class with \(p\nmid qa\), and write
\(I'\mathfrak p'=(y)\). Apply the principal calculation to
\(y\) and the prime calculation to \(\mathfrak p'\). Divide
the two congruences, using units at each place, and apply uniqueness
of the choices at both levels. The resulting root \(\eta\) is a
product of the same globally chosen roots at every place, as required.

\subsection{The family at the fixed prime}

Let \(\mathfrak r_a\) be the prime cut out on \(F_a\) by our
fixed place at \(r\). For every admissible product \(n\) take
\(X_n=X_{Pn}(\mathfrak r_{Pn})\); arbitrary choices suffice.
Since \(\mathfrak r_{a}\) lies above \(r\), their divisors show
that they are units away from \(r\). For \(a=P\), \(r\equiv1\bmod P\)
implies complete splitting and unramifiedness. Thus for every
\(t\in G_P\), the valuation of \(\sigma_t X_1\) at the fixed
place, normalized by \(v(r)=1\), is precisely \(h_P(t)\).

At the layer \(a\mapsto qa\) put
\(s=[k(\mathfrak r_{qa}):k(\mathfrak r_a)]\).
The norm of \(\mathfrak r_{qa}\) is \(\mathfrak r_a^s\), for
which \(X_n^s\) is a normalized generator. The integer \(s\)
divides the absolute residue degree
\(\operatorname{ord}_{qa}(r)\), since \(r\nmid qa\).
By the Chinese remainder theorem, the latter divides the product
of the orders of \(r\bmod p\) over primes \(p\mid qn\), as
\(r\equiv1\bmod P\) and those primes are distinct and coprime
to \(P\). Every factor is prime to \(\ell\) by admissibility,
so \(\ell\nmid s\). Applying the norm and simultaneous reduction
just proved to \(X_{qn}\) and \(X_n^s\) gives exactly the two input
relations in Section~\ref{sec:minus-derivatives}.
\footnote{\href{https://github.com/ChristianKudeba/alpoge-rh-verification/blob/alpoge-source-7cbf7a0c/src/Development/AutomorphicArithmetic/AutomorphicVarieties/PrimitiveFrobeniusFixed/XJacPowNe.lean}{Jacobi sums and their prime divisors},
\href{https://github.com/ChristianKudeba/alpoge-rh-verification/blob/alpoge-source-7cbf7a0c/src/Development/AutomorphicArithmetic/AutomorphicVarieties/PrimitiveFrobeniusFixed/XEx.lean}{passage to arbitrary ideals},
\href{https://github.com/ChristianKudeba/alpoge-rh-verification/blob/alpoge-source-7cbf7a0c/src/Development/AutomorphicArithmetic/AutomorphicVarieties/PrimitiveFrobeniusFixed/ThDist.lean}{distribution of the signed floor coefficients},
\href{https://github.com/ChristianKudeba/alpoge-rh-verification/blob/alpoge-source-7cbf7a0c/src/Development/AutomorphicArithmetic/AutomorphicVarieties/PrimitiveFrobeniusFixed/XNormDiv.lean}{norm relation},
\href{https://github.com/ChristianKudeba/alpoge-rh-verification/blob/alpoge-source-7cbf7a0c/src/Development/AutomorphicArithmetic/AutomorphicVarieties/PrimitiveFrobeniusFixed/XCongJac.lean}{Jacobi congruence},
\href{https://github.com/ChristianKudeba/alpoge-rh-verification/blob/alpoge-source-7cbf7a0c/src/Development/AutomorphicArithmetic/AutomorphicVarieties/PrimitiveFrobeniusFixed/XCongPrin.lean}{principal congruence},
\href{https://github.com/ChristianKudeba/alpoge-rh-verification/blob/alpoge-source-7cbf7a0c/src/Development/AutomorphicArithmetic/AutomorphicVarieties/PrimitiveFrobeniusFixed/AlgStep.lean}{division to obtain simultaneous congruences},
\href{https://github.com/ChristianKudeba/alpoge-rh-verification/blob/alpoge-source-7cbf7a0c/src/Development/AutomorphicArithmetic/AutomorphicVarieties/PrimitiveFrobeniusFixed/PlaceEmb.lean}{assembly at the chosen place}.}
